\documentclass[10pt,a4paper]{article}

\usepackage[margin=0.58in]{geometry}
\usepackage[T1]{fontenc}
\usepackage{lmodern}
\usepackage{microtype}
\usepackage{enumitem}
\usepackage{tabularx}
\usepackage{array}
\usepackage{titlesec}
\usepackage[hidelinks]{hyperref}
\usepackage{xcolor}
\usepackage{ragged2e}
\usepackage{setspace}

\pagestyle{empty}
\setlength{\parindent}{0pt}
\setlength{\parskip}{0pt}
\setlength{\tabcolsep}{0pt}
\renewcommand{\arraystretch}{1.0}
\setstretch{1.0}

\definecolor{cvgray}{HTML}{555555}

\titleformat{\section}
  {\large\bfseries}
  {}{0pt}{}
  [\vspace{1pt}\titlerule]
\titlespacing*{\section}{0pt}{7pt}{4pt}

\newcommand{\cvitem}[2]{%
  \begin{tabularx}{\textwidth}{@{}X >{\raggedleft\arraybackslash}p{0.23\textwidth}@{}}
    #1 & #2
  \end{tabularx}%
}

\newcommand{\cventry}[4]{%
  \cvitem{\textbf{#1}}{\textcolor{cvgray}{#2}}\\[-1pt]
  {#3}\\[-1pt]
  {\small\textcolor{cvgray}{#4}}\vspace{2pt}
}

\newcommand{\researchentry}[4]{%
  \begin{tabularx}{\textwidth}{@{}X >{\raggedleft\arraybackslash}p{0.10\textwidth}@{}}
    \textbf{#1} & \textcolor{cvgray}{#2}
  \end{tabularx}\\[-1pt]
  {\small #3}\vspace{1pt}
  \begin{itemize}[leftmargin=1.25em,itemsep=1pt,topsep=1pt,parsep=0pt]
    #4
  \end{itemize}
  \vspace{0pt}
}

\newcommand{\ongoingentry}[3]{%
  \begin{tabularx}{\textwidth}{@{}X >{\raggedleft\arraybackslash}p{0.10\textwidth}@{}}
    \textbf{#1} & \textcolor{cvgray}{#2}
  \end{tabularx}
  \begin{itemize}[leftmargin=1.25em,itemsep=1pt,topsep=1pt,parsep=0pt]
    #3
  \end{itemize}
  \vspace{0pt}
}

\newcommand{\linkrow}[2]{%
  {\small #1 \;|\; #2}\vspace{2pt}
}

\newcommand{\researchsub}[1]{%
  \vspace{2pt}{\small\bfseries\underline{#1}}\par\vspace{1pt}
}

\begin{document}

% ===== Header =====
{\LARGE\bfseries SEONVIN CHO}\\[3pt]
Integrated M.S./Ph.D. Student in Electronic Engineering, Hanyang University\\[2pt]
\href{mailto:seonbin0319@hanyang.ac.kr}{seonbin0319@hanyang.ac.kr}
\;|\;
\href{https://seonvin0319.github.io/index.html}{Website}
\;|\;
\href{https://github.com/seonvin0319}{GitHub}
\;|\;
Seoul, Republic of Korea

\section*{Research Interests}
Reinforcement Learning \;\textperiodcentered\; Sequential Decision Making \;\textperiodcentered\; Generative Models

\section*{Education}
\cventry
  {Hanyang University, Seoul, Republic of Korea}
  {Sep. 2024 -- Present}
  {Integrated M.S./Ph.D. Student in Electronic Engineering}
  {\textit{Advisor: \href{https://sites.google.com/view/snlab}{Prof. Songnam Hong}}}\\[-1pt]
\textbf{GPA:} 4.45 / 4.50
\vspace{3pt}

\cventry
  {Hanyang University, Seoul, Republic of Korea}
  {Mar. 2020 -- Aug. 2024}
  {B.S. in Electronic Engineering}
  {\textit{Summa Cum Laude, Early Graduation}}\\[-1pt]
\textbf{GPA:} 4.26 / 4.50

\section*{Research}

\researchsub{Recent Work}
\researchentry
  {PathBridger: Subgoal Bridges for Offline Goal-Conditioned Reinforcement Learning}
  {2026}
  {Soohyun Choi$^{*}$, \textbf{Seonvin Cho$^{*}$}, \href{https://sites.google.com/view/snlab}{Songnam Hong}$^{\dagger}$ \;\textperiodcentered\; \textit{Preprint} \;\textperiodcentered\; $^{*}$Equal contribution \;\textperiodcentered\; $^{\dagger}$Corresponding author\\
   \href{https://github.com/SChoish/PathBridger}{[Code]} \;\; \href{https://arxiv.org/pdf/2608.29061}{[arXiv]}}
  {
    \item Developed a hierarchical offline goal-conditioned RL framework that connects subgoal selection to short-horizon execution through explicit state-space bridges.
    \item Combined generative endpoint proposals, transitive value-based selection, and inverse-dynamics decoding to produce executable short-horizon action chunks.
  }

\researchentry
  {Multi-step Proximal Policy Improvement in Offline Reinforcement Learning}
  {2026}
  {Soohyun Choi$^{*}$, \textbf{Seonvin Cho$^{*}$}, \href{https://sites.google.com/view/snlab}{Songnam Hong}$^{\dagger}$ \;\textperiodcentered\; \textit{Preprint} \;\textperiodcentered\; $^{*}$Equal contribution \;\textperiodcentered\; $^{\dagger}$Corresponding author\\
   \href{https://github.com/SChoish/MPI}{[Code]} \;\; \href{https://arxiv.org/pdf/2609.03842}{[arXiv]}}
  {
    \item Developed a geometric view of offline actor updates, interpreting a broad class of behavior-anchored objectives as a single proximal policy improvement step on a statistical manifold.
    \item Proposed MPI, a plug-in refinement that composes sequential re-centered proximal steps and improves TD3+BC, ReBRAC, and IQL on many D4RL tasks.
  }

\researchentry
  {Federated Learning under Noisy Labels}
  {2023}
  {Undergraduate Capstone Project\\
   \href{https://github.com/seonvin0319/TriMix}{[Code]} \;\; \href{https://seonvin0319.github.io/assets/Analyzing_Loss_Value_Distribution_from_Learning_Label_Noise_Data_in_the_Federated_Learning_Environment.pdf}{[Technical Report]}}
  {
    \item Analyzed loss-value distributions in federated learning to identify unreliable samples under label noise.
    \item Received the Best Award (Department) and Excellence Award (College of Engineering).
  }

\researchsub{Ongoing Research}
\ongoingentry
  {MART: Multi-Step Actor Refinement Trajectories}
  {Present}
  {
    \item Developing a fixed-budget actor-refinement path that divides a nominal improvement horizon across predecessor-centered actors and separates critic-facing exposure from deployment reach.
  }

\ongoingentry
  {AMO: Adaptive Multiscale Policy Optimization}
  {Present}
  {
    \item Studying adaptive proximal policy improvement that learns a total horizon \(T\) with a cross-fitted outer objective, using a conservative \(T/N\)-policy for Bellman targets and a full-\(T\) policy for execution.
  }

\ongoingentry
  {Extensions of PathBridger}
  {Present}
  {
    \item Exploring offline-to-online learning, longer-horizon control, and action-free settings where the data do not provide actions.
  }

\section*{Honors \& Awards}
\cvitem{BK Research Assistantship (BK-RA), Hanyang University}{Spring \& Fall 2026}\\[2pt]
\cvitem{IEEE Antennas and Propagation Society Undergraduate Summer Research Scholarship}{2024}\\[2pt]
\cvitem{Excellence Paper Award, KIEES \;\; \href{https://github.com/seonvin0319/gpu_dispersive_fdtd_ss}{[Code]}}{Winter \& Fall 2024}\\[2pt]
\cvitem{Best \& Excellence Awards, Federated Learning Capstone Project}{2023}\\[2pt]
\cvitem{Excellence Award, Embedded Creative Robot Challenge}{2022}\\[2pt]
\cvitem{Merit-Based \& Academic Scholarships, Hanyang University}{2021 -- 2024}

\end{document}
